% GENERATED FILE. Edit canonical lessons, then run npm run generate:books.
% canonical-source-hash: fnv1a64:1a1b42776542ea9b
% canonical-lessons: JA-C30-osu, JA-C30-hiku, JA-C30-sagasu, JA-C30-kasu, JA-C30-kaesu

\chapter{To push, To pull, To look for, To lend, To give back}
\label{ch:ja-to-push-to-pull-to-look-for-to-lend-to-give-back}
\hlchaptermodality{30}

\begin{chapteropening}
\textbf{By the end of this chapter:} \emph{I can say to push, to pull, to look for, to lend and to give back.}

Complete the last lesson of chapter 30: i can say to push, to pull, to look for, to lend and to give back.
\end{chapteropening}

\section[osu]{\ja{おす} (osu) — to push}

\label{lesson:JA-C30-osu}



\begin{quote}
Before the new one: say the Japanese for to rest, then the Japanese for to enjoy.
\end{quote}

\subsection*{You'll want to know: \ja{おす}}
\textbf{\ja{おす}} — \emph{osu} — \textquotedblleft{}to push\textquotedblright{}.

\begin{grammarlens}[title={What you've built}]
One more word for this chapter.
\end{grammarlens}

\subsection*{Your turn}
\begin{itemize}
\raggedright
  \item \emph{Say it:} \emph{osu}
  \item \emph{Say it:} \emph{osu}, once more
  \item \emph{Say it:} say \emph{osu}
  \item \emph{Recall:} say the Japanese for to rest, then the Japanese for to enjoy, then say \emph{osu} again
\end{itemize}

\begin{tcolorbox}[breakable,colback=teal!4,colframe=teal!35!black,title={Before you move on}]
What does \ja{おす} mean? (\textquotedblleft{}To push\textquotedblright{}.) Say it once more.
\end{tcolorbox}
\section[hiku]{\ja{ひく} (hiku) — to pull}

\label{lesson:JA-C30-hiku}



\begin{quote}
Before the new one: say the Japanese for to enjoy, then the Japanese for to push.
\end{quote}

\subsection*{You'll want to know: \ja{ひく}}
\textbf{\ja{ひく}} — \emph{hiku} — \textquotedblleft{}to pull\textquotedblright{}.

\begin{grammarlens}[title={What you've built}]
One more word for this chapter.
\end{grammarlens}

\subsection*{Your turn}
\begin{itemize}
\raggedright
  \item \emph{Say it:} \emph{hiku}
  \item \emph{Say it:} \emph{hiku}, once more
  \item \emph{Say it:} say \emph{hiku}
  \item \emph{Recall:} say the Japanese for to enjoy, then the Japanese for to push, then say \emph{hiku} again
\end{itemize}

\begin{tcolorbox}[breakable,colback=teal!4,colframe=teal!35!black,title={Before you move on}]
What does \ja{ひく} mean? (\textquotedblleft{}To pull\textquotedblright{}.) Say it once more.
\end{tcolorbox}
\section[sagasu]{\ja{さがす} (sagasu) — to look for}

\label{lesson:JA-C30-sagasu}



\begin{quote}
Before the new one: say the Japanese for to push, then the Japanese for to pull.
\end{quote}

\subsection*{You'll want to know: \ja{さがす}}
\textbf{\ja{さがす}} — \emph{sagasu} — \textquotedblleft{}to look for\textquotedblright{}.

\begin{grammarlens}[title={What you've built}]
One more word for this chapter.
\end{grammarlens}

\subsection*{Your turn}
\begin{itemize}
\raggedright
  \item \emph{Say it:} \emph{sagasu}
  \item \emph{Say it:} \emph{sagasu}, once more
  \item \emph{Say it:} say \emph{sagasu}
  \item \emph{Recall:} say the Japanese for to push, then the Japanese for to pull, then say \emph{sagasu} again
\end{itemize}

\begin{tcolorbox}[breakable,colback=teal!4,colframe=teal!35!black,title={Before you move on}]
What does \ja{さがす} mean? (\textquotedblleft{}To look for\textquotedblright{}.) Say it once more.
\end{tcolorbox}
\section[kasu]{\ja{かす} (kasu) — to lend}

\label{lesson:JA-C30-kasu}



\begin{quote}
Before the new one: say the Japanese for to pull, then the Japanese for to look for.
\end{quote}

\subsection*{You'll want to know: \ja{かす}}
\textbf{\ja{かす}} — \emph{kasu} — \textquotedblleft{}to lend\textquotedblright{}.

\begin{grammarlens}[title={What you've built}]
One more word for this chapter.
\end{grammarlens}

\subsection*{Your turn}
\begin{itemize}
\raggedright
  \item \emph{Say it:} \emph{kasu}
  \item \emph{Say it:} \emph{kasu}, once more
  \item \emph{Say it:} say \emph{kasu}
  \item \emph{Recall:} say the Japanese for to pull, then the Japanese for to look for, then say \emph{kasu} again
\end{itemize}

\begin{tcolorbox}[breakable,colback=teal!4,colframe=teal!35!black,title={Before you move on}]
What does \ja{かす} mean? (\textquotedblleft{}To lend\textquotedblright{}.) Say it once more.
\end{tcolorbox}
\section[kaesu]{\ja{かえす} (kaesu) — to give back}

\label{lesson:JA-C30-kaesu}



\begin{quote}
Before the new one: say the Japanese for to look for, then the Japanese for to lend.
\end{quote}

\subsection*{You'll want to know: \ja{かえす}}
\textbf{\ja{かえす}} — \emph{kaesu} — \textquotedblleft{}to give back\textquotedblright{}.

\begin{grammarlens}[title={What you've built}]
One more word for this chapter.
\end{grammarlens}

\subsection*{Your turn}
\begin{itemize}
\raggedright
  \item \emph{Say it:} \emph{kaesu}
  \item \emph{Say it:} \emph{kaesu}, once more
  \item \emph{Say it:} say \emph{kaesu}
  \item \emph{Recall:} say the Japanese for to look for, then the Japanese for to lend, then say \emph{kaesu} again
\end{itemize}

\begin{tcolorbox}[breakable,colback=teal!4,colframe=teal!35!black,title={Before you move on}]
What does \ja{かえす} mean? (\textquotedblleft{}To give back\textquotedblright{}.) Say it once more.
\end{tcolorbox}
